\documentclass[11pt]{article}
\usepackage[a4paper,margin=26mm]{geometry}
\usepackage{amsmath,amssymb,amsthm,mathtools}
\usepackage[colorlinks=true,urlcolor=blue,citecolor=blue]{hyperref}
\newtheorem{theorem}{Theorem}
\newtheorem{lemma}[theorem]{Lemma}
\newtheorem{corollary}[theorem]{Corollary}
\newtheorem{proposition}[theorem]{Proposition}
\newcommand{\Ric}{\operatorname{Ric}}
\newcommand{\R}{\mathbb R}
\newcommand{\Z}{\mathbb Z}
\newcommand{\CP}{\mathbb{CP}}
\newcommand{\RP}{\mathbb{RP}}
\newcommand{\Vol}{\operatorname{Vol}}
\title{A Global Topological Exclusion of Ricci-Flat ALE Cores\\
\large A continuation toward smooth curvature-energy minimizers}
\author{Research continuation}
\date{11 October 2026}
\begin{document}
\maketitle
\begin{abstract}
We derive an obstruction to embedding a nonflat Ricci-flat ALE four-manifold's
compact core in a closed aspherical four-manifold. No self-duality or
classification of general Ricci-flat ALE metrics is assumed. More generally,
a finite-fundamental-group ALE core with nontrivial group at
infinity has traceless-Ricci energy at least $4\pi^2+\frac1{12}\int s^2$ if it embeds in such an
ambient manifold. The constant is attained by a classical-metric quotient.
The proof combines integral topology, the classification of spherical
three-manifold groups, and ALE characteristic formulas. These results are
consequences of the cited classical inputs; no priority claim is made.
They remove the zero-energy ALE channel for a different class of candidate
manifolds, but do not prove existence of a smooth non-Bach-flat global
curvature-energy minimizer.
\end{abstract}

\section{Statement and conventions}
All manifolds are smooth and oriented. A closed manifold is aspherical when its
universal cover is contractible. Put
\[
 P(g)=\int |\Ric_g^{\circ}|^2\,dV_g,\qquad
 W_\pm(g)=\int|W_g^\pm|^2\,dV_g.
\]
Weyl norms are curvature-operator norms; the Ricci norm is the usual tensor
norm. An ALE manifold here is complete, connected, one-ended, and has
an end $(\R^4\setminus B)/\Gamma$, with the standard differentiated positive-order
ALE decay. A compact core is obtained by cutting off this end sufficiently far
out. Embedding a core always means a smooth codimension-zero embedding; no
metric matching is part of this terminology.

\begin{theorem}\label{thm:main}
Let $(N,g)$ be an ALE four-manifold with finite fundamental group.
Suppose a compact core of $N$ embeds in a closed aspherical four-manifold $M$.
If $\Gamma\ne1$, then
\[
                         P(g)\ge4\pi^2+\frac1{12}\int_Ns_g^2\,dV_g.
\]
Consequently, no compact core of a nonflat Ricci-flat ALE four-manifold embeds
in a closed aspherical four-manifold.
\end{theorem}

The second conclusion includes Ricci-flat metrics with either Weyl component
nonzero. Its proof below does not infer self-duality from Ricci-flatness.

\section{Topology of the embedded core}
\begin{lemma}\label{lem:topology}
Let $N_0\subset M$ be a compact connected codimension-zero submanifold, with
connected spherical-space-form boundary $Y$, and with finite $\pi_1(N_0)$.
If $M$ is closed and aspherical, then
\[
 H_2(N_0;\Z)=0,\quad \chi(N_0)=1,\quad \tau(N_0)=0,
 \qquad H_1(Y;\Z)\cong A\oplus A,
 \quad A=H_1(N_0;\Z).
\]
\end{lemma}
\begin{proof}
The group of a finite-dimensional aspherical manifold is torsion-free:
its integral cohomological dimension is finite, whereas a group containing
nontrivial torsion has infinite integral cohomological dimension. The map
$\pi_1(N_0)\to\pi_1(M)$ is therefore trivial. The embedding lifts to an
embedding of $N_0$ in the contractible universal cover $U$.

Write $C=\overline{U\setminus N_0}$, using collars for Mayer--Vietoris.
The boundary is connected and the exterior is connected.
Since $H_2(Y;\Z)=0$ and $U$ is contractible, Mayer--Vietoris gives
\[
 H_2(N_0;\Z)=H_2(C;\Z)=0,\qquad
 H_1(Y;\Z)\xrightarrow{\ \cong\ } A\oplus H_1(C;\Z).
\]
The fundamental-class boundary map
$H_4(N_0,Y;\Z)\to H_3(Y;\Z)$ is an isomorphism.
Moreover, $H_3(N_0,Y;\Z)\cong H^1(N_0;\Z)=0$, because $A$ is finite.
Thus $H_3(N_0;\Z)=0$.

Poincar\'e--Lefschetz duality and the universal coefficient theorem give
\[
 H_2(N_0,Y;\Z)\cong H^2(N_0;\Z)\cong\operatorname{Ext}(A,\Z),
 \qquad H_1(N_0,Y;\Z)\cong H^3(N_0;\Z)=0.
\]
The pair sequence is consequently
\[
 0\longrightarrow\operatorname{Ext}(A,\Z)
 \longrightarrow H_1(Y;\Z)\longrightarrow A\longrightarrow0.
\]
In the Mayer--Vietoris splitting the last map is projection onto $A$.
Its kernel is therefore $H_1(C;\Z)\cong\operatorname{Ext}(A,\Z)\cong A$.
This proves the direct-double assertion. The remaining real homology
calculation gives $\chi=1$ and $\tau=0$.
\end{proof}

For instance, a standard nontrivial cyclic rational ball has
$H_1(N_0)=\Z/d$ and $H_1(Y)=\Z/d^2$, with $d>1$; these cannot form the
required direct double. This argument uses contractibility of the universal
cover, not merely torsion-freeness of the ambient fundamental group.
It does not apply to a blowup of an aspherical manifold.

\section{Boundary groups and the curvature inequality}
The group-theoretic part of the proof of Hillman's Theorem 5.2
\cite{Hillman} says that a finite spherical three-manifold group whose
abelianization is a direct double is one of
\[
                 1,\qquad I^*,\qquad Q(8k)\quad(k\ge1).
\]
Here $I^*$ is the binary icosahedral group of order $120$, and $Q(8k)$
is the generalized quaternionic group of order $8k$.
We use this algebraic restriction only. We do not invoke the stronger theorem
about embedding in a compact homology four-sphere, and do not compactify $U$.

Use the signature-eta convention
\begin{align}
 8\pi^2\left(\chi(N)-\frac1{|\Gamma|}\right)
   &=W_++W_- -\frac12P+\frac1{24}\int_Ns^2,\label{eq:gb}\\
 \tau(N)&=\frac{W_+-W_-}{12\pi^2}-\eta(S^3/\Gamma).
 \label{eq:sig}
\end{align}
The ALE end supplies the displayed boundary correction.
Equations \eqref{eq:gb}--\eqref{eq:sig}, with
Lemma~\ref{lem:topology}, imply
\begin{equation}\label{eq:cost}
 P\ge24\pi^2|\eta(S^3/\Gamma)|
            -16\pi^2\left(1-\frac1{|\Gamma|}\right)+\frac1{12}\int_Ns^2.
\end{equation}
This uses only $W_++W_-\ge|W_+-W_-|$.

For the quaternionic groups, Wright's eta calculation
\cite[Lemma 1]{Wright}, with $u=1$ and $v=2k$, gives, up to orientation,
\[
 |\eta(S^3/Q(8k))|=\frac{2k}{3}+\frac1{12k}.
\]
Therefore
\begin{equation}\label{eq:quaternion}
 P\ge16\pi^2(k-1)+\frac{4\pi^2}{k}+\frac1{12}\int_Ns^2
 \ge4\pi^2+\frac1{12}\int_Ns^2.
\end{equation}
For $I^*$, its eta invariant can equivalently be computed using the known
$E_8$ hyperk\"ahler ALE resolution, with $\chi=9$, $\tau=-8$ and end order
$120$. The same formulas give
\[
 |\eta(S^3/I^*)|=8-\frac23\left(9-\frac1{120}\right)
               =\frac{361}{180}.
\]
Hence \eqref{eq:cost} becomes
\[
                              P\ge\frac{484}{15}\pi^2+\frac1{12}\int_Ns^2.
\]
This proves the first assertion of Theorem~\ref{thm:main}.

For the Ricci-flat conclusion, the fundamental group is finite by the standard
Bishop--Gromov covering argument: nonnegative Ricci curvature and positive
asymptotic volume ratio bound the number of deck transformations. More
explicitly, disjoint translated fundamental domains associated with any fixed
finite collection of deck transformations give, as the radius tends to
infinity, that collection's cardinality times the base asymptotic volume ratio
as a lower bound for the covering ball's volume ratio. Bishop--Gromov bounds
the latter by one. In the ALE case the base ratio is $1/|\Gamma|$.

If $\Gamma\ne1$, the positive lower bound contradicts $P=0$.
If $\Gamma=1$, \eqref{eq:gb} and $P=0$, $\chi=1$ imply
$W_++W_-=0$. All curvature then vanishes. This completes the proof.

\section{Sharpness and the remaining positive-energy channel}
The constant $4\pi^2$ is attained within the class of cores just considered.
Let $g_4$ be LeBrun's radial scalar-flat K\"ahler ALE metric on
$\mathcal O(-4)$ \cite{LeBrun88,Honda}. On its end, let
\[
              j(z_1,z_2)=(-\overline z_2,\overline z_1).
\]
The end is $\mathbb C^2/\Z_4$ with generator $i\operatorname{Id}$.
The map $j$ normalizes this action and satisfies $j^2=-\operatorname{Id}$.
It extends to a free isometric involution of $\mathcal O(-4)$: on the zero
section it is antipodal, and away from it the equation $jz=i^a z$ has no
nonzero solution. In complex dimension two it preserves the real orientation.
The quotient $N$ is smooth, scalar-flat, anti-self-dual and ALE, with end
$S^3/Q(8)$.

Its core is the disk bundle over $\RP^2$ of normal Euler number of absolute
value two, the neighborhood of the standard embedded $\RP^2\subset S^4$
\cite{Hillman,KS}. Removing a point outside this neighborhood embeds the
core in $\R^4$, and hence in a coordinate ball of any four-manifold.

For the scalar-flat K\"ahler metric $g_m$ on $\mathcal O(-m)$, the
anti-self-dual Ricci form and the compactly supported intersection pairing give
\[
 c_1^2=-\frac{(m-2)^2}{m},\qquad
 P(g_m)=-8\pi^2c_1^2=8\pi^2\frac{(m-2)^2}{m}.
\]
At $m=4$ the energy is $8\pi^2$; division by the free involution yields
\[
                                P(N)=4\pi^2.
\]
Thus asphericity cannot eliminate all scalar-flat ALE cores. This is an
explicit assembly of classical geometry and topology, not a claim that such
a core occurs in a global minimizing sequence. Nor does its topological
embedding in a ball by itself justify a lower-energy metric replacement:
the controlled ALE boundary is $S^3/Q(8)$, while the outer $S^3$ of the ball
has an additional intervening region whose geometry must be controlled.

\section{Connection with the closed minimizer problem}
Capovilla constructs closed oriented aspherical smooth four-manifolds $X_n$
with $\chi(X_n)=\tau(X_n)=n>0$ \cite[Theorem 1]{Capovilla}.
For every smooth metric on $X_n$, the opposite-chirality Chern--Weil identity is
\[
                  P(g)=8\pi^2n+4W_-(g)+\frac1{12}\int s_g^2.
\]
In particular, Einstein metrics are excluded. Chodosh--Li's theorem says that
a nonnegative-scalar-curvature metric on a closed aspherical four-manifold is
flat \cite[Theorem 2]{CL}. Since $\chi(X_n)>0$, every conformal class here has
negative Yamabe constant. This does not assert that the supremum of these
constants is strictly negative.

Any smooth $P$-critical metric on $X_n$ has constant scalar curvature by
conformal variation, and that constant must therefore be negative.
The closed identity
\[
                P=2(W_++W_-)+\frac1{12}\int s^2-16\pi^2\chi
\]
shows that a Bach-flat $P$-critical metric of nonzero constant scalar curvature
is Einstein: the Weyl gradient vanishes, leaving the scalar-square gradient
$-2s\Ric^\circ$ equal to zero. Consequently a smooth global $P$-minimizer on
one $X_n$ would answer LeBrun's non-Bach-flat optimal-metric question
\cite[Section 11]{LeBrun04}.

\medskip\noindent\textbf{Exact remaining issue.}
No such smooth global minimizer is proved here. Anderson's ideal-attainment
theorem \cite[Theorem 4.10 and Remark 4.12]{Anderson} still permits singular
configurations and an open or partially collapsed branch. Theorem~\ref{thm:main}
removes all nonflat Ricci-flat ALE cores from an embedded-core decomposition
on these candidates. Positive-$P$ cores remain, including the sharp example
above; general scalar-flat ALE cores are not here proved to have finite
fundamental group; and partial collapse remains uncontrolled. A global
comparison or compactness argument addressing those possibilities is still
required. Changing the candidate manifold preserves the unrestricted
functional and the original existence question; it does not establish its
answer.

\begin{thebibliography}{9}
\bibitem{Hillman} J. A. Hillman, \emph{Locally flat embeddings of 3-manifolds
in $S^4$}, arXiv:2408.10535, proof of Theorem 5.2.
\url{https://arxiv.org/abs/2408.10535}
\bibitem{Wright} E. P. Wright, \emph{Quotients of gravitational instantons},
arXiv:1102.2442, Lemma 1 and Sections 2, 6, 8.
\url{https://arxiv.org/abs/1102.2442}
\bibitem{LeBrun88} C. LeBrun, \emph{Counter-examples to the generalized
positive action conjecture}, Comm. Math. Phys. 118 (1988), 591--596.
\bibitem{Honda} N. Honda, \emph{Deformation of LeBrun's ALE metrics with
negative mass}, arXiv:1204.4857.
\url{https://arxiv.org/abs/1204.4857}
\bibitem{KS} A. Katanaga and O. Saeki, \emph{Embeddings of quaternion space
in $S^4$}, J. Austral. Math. Soc. 65 (1998), 313--325.
\bibitem{Capovilla} P. Capovilla, \emph{Aspherical 4-manifolds with positive
Euler characteristic and their geography}, arXiv:2511.15577.
\url{https://arxiv.org/abs/2511.15577}
\bibitem{CL} O. Chodosh and C. Li, \emph{Generalized soap bubbles and the
topology of manifolds with positive scalar curvature}, Ann. of Math. 199
(2024), 707--740. \url{https://arxiv.org/abs/2008.11888}
\bibitem{LeBrun04} C. LeBrun, \emph{Curvature functionals, optimal metrics,
and the differential topology of 4-manifolds}, arXiv:math/0404251.
\url{https://arxiv.org/abs/math/0404251}
\bibitem{Anderson} M. T. Anderson, \emph{Canonical metrics on 3-manifolds
and 4-manifolds}, arXiv:math/0511034.
\url{https://arxiv.org/abs/math/0511034}
\end{thebibliography}
\end{document}
